% GENERATED FILE. Edit canonical lessons, then run npm run generate:books.
% canonical-source-hash: fnv1a64:a6f8af73d949f2b2
% canonical-lessons: ES-C659-lancha, ES-C659-velero, ES-C659-remar, ES-C659-marinero, ES-C659-capitan

\chapter{Barcas y marineros}
\label{ch:es-barcas-y-marineros}
\hlchaptermodality{659}

\begin{chapteropening}
\textbf{By the end of this chapter:} \emph{I can talk about small boats: a launch, a sailing boat, rowing, the sailor and the captain.}

Complete the last lesson of chapter 659: I can talk about small boats: a launch, a sailing boat, rowing, the sailor and the captain.
\end{chapteropening}

\section[la lancha]{la lancha — the launch, the motorboat}

\label{lesson:ES-C659-lancha}



\begin{quote}
Before the new one: say the Spanish for the quay, then the Spanish for to disembark.
\end{quote}

\subsection*{You'll want to know: la lancha}
\textbf{la lancha} — \textquotedblleft{}the launch, the motorboat\textquotedblright{}.

\begin{quote}
Fuimos a la cueva en lancha.
\end{quote}

\begin{quote}
\emph{We went to the cave by motorboat.}
\end{quote}

\emph{La lancha} has a motor; \emph{la barca} is often rowed.

\begin{cousinweb}
\emph{lancha} came through Portuguese \emph{lancha} and is usually traced to a Malay word meaning \textquotedblleft{}swift\textquotedblright{}.
\end{cousinweb}

\begin{grammarlens}[title={What you've built}]
One more word for telling, explaining and describing at B1.
\end{grammarlens}

\subsection*{Your turn}
\begin{itemize}
\raggedright
  \item \emph{Say it:} \emph{la lancha}
  \item \emph{Say it:} \emph{la lancha}, once more
  \item \emph{Say it:} a sentence of your own with \emph{la lancha}
  \item \emph{Recall:} say the Spanish for the quay, then the Spanish for to disembark, then say \emph{la lancha} again
\end{itemize}

\begin{tcolorbox}[breakable,colback=teal!4,colframe=teal!35!black,title={Before you move on}]
What does la lancha mean? (\textquotedblleft{}The launch, the motorboat\textquotedblright{}.) Say it once more.
\end{tcolorbox}
\section[el velero]{el velero — the sailing boat}

\label{lesson:ES-C659-velero}



\begin{quote}
Before the new one: say the Spanish for to disembark, then the Spanish for the launch.
\end{quote}

\subsection*{You'll want to know: el velero}
\textbf{el velero} — \textquotedblleft{}the sailing boat\textquotedblright{}.

\begin{quote}
El velero sale del puerto con poco viento.
\end{quote}

\begin{quote}
\emph{The sailing boat leaves the harbour in a light wind.}
\end{quote}

Built on \emph{la vela}, the sail.

\begin{cousinweb}
\emph{velero} is built on \emph{vela} \textquotedblleft{}sail\textquotedblright{}, which comes from Latin \emph{vēla}, the plural of \emph{vēlum} \textquotedblleft{}a cloth, a covering\textquotedblright{} (Lewis \& Short), used for sails.
\end{cousinweb}

\begin{grammarlens}[title={What you've built}]
One more word for telling, explaining and describing at B1.
\end{grammarlens}

\subsection*{Your turn}
\begin{itemize}
\raggedright
  \item \emph{Say it:} \emph{el velero}
  \item \emph{Say it:} \emph{el velero}, once more
  \item \emph{Say it:} a sentence of your own with \emph{el velero}
  \item \emph{Recall:} say the Spanish for to disembark, then the Spanish for the launch, then say \emph{el velero} again
\end{itemize}

\begin{tcolorbox}[breakable,colback=teal!4,colframe=teal!35!black,title={Before you move on}]
What does el velero mean? (\textquotedblleft{}The sailing boat\textquotedblright{}.) Say it once more.
\end{tcolorbox}
\section[remar]{remar — to row}

\label{lesson:ES-C659-remar}



\begin{quote}
Before the new one: say the Spanish for the launch, then the Spanish for the sailing boat.
\end{quote}

\subsection*{You'll want to know: remar}
\textbf{remar} — \textquotedblleft{}to row\textquotedblright{}.

\begin{quote}
Remamos hasta la isla.
\end{quote}

\begin{quote}
\emph{We rowed to the island.}
\end{quote}

\emph{El remo} is the oar.

\begin{cousinweb}
\emph{remar} is built on \emph{remo} \textquotedblleft{}oar\textquotedblright{} (EtymDB), inherited from Latin \emph{rēmus} \textquotedblleft{}an oar\textquotedblright{} (Lewis \& Short).
\end{cousinweb}

\begin{grammarlens}[title={What you've built}]
One more word for telling, explaining and describing at B1.
\end{grammarlens}

\subsection*{Your turn}
\begin{itemize}
\raggedright
  \item \emph{Say it:} \emph{remar}
  \item \emph{Say it:} \emph{remar}, once more
  \item \emph{Say it:} a sentence of your own with \emph{remar}
  \item \emph{Recall:} say the Spanish for the launch, then the Spanish for the sailing boat, then say \emph{remar} again
\end{itemize}

\begin{tcolorbox}[breakable,colback=teal!4,colframe=teal!35!black,title={Before you move on}]
What does remar mean? (\textquotedblleft{}To row\textquotedblright{}.) Say it once more.
\end{tcolorbox}
\section[el marinero]{el marinero — the sailor}

\label{lesson:ES-C659-marinero}



\begin{quote}
Before the new one: say the Spanish for the sailing boat, then the Spanish for to row.
\end{quote}

\subsection*{You'll want to know: el marinero}
\textbf{el marinero} — \textquotedblleft{}the sailor\textquotedblright{}.

\begin{quote}
Mi abuelo era marinero.
\end{quote}

\begin{quote}
\emph{My grandfather was a sailor.}
\end{quote}

As an adjective it means seafaring: \emph{un pueblo marinero}.

\begin{cousinweb}
\emph{marinero} is built on \emph{marina} (EtymDB), and through it on \emph{mar} \textquotedblleft{}sea\textquotedblright{}, inherited from Latin \emph{mare}.
\end{cousinweb}

\begin{grammarlens}[title={What you've built}]
One more word for telling, explaining and describing at B1.
\end{grammarlens}

\subsection*{Your turn}
\begin{itemize}
\raggedright
  \item \emph{Say it:} \emph{el marinero}
  \item \emph{Say it:} \emph{el marinero}, once more
  \item \emph{Say it:} a sentence of your own with \emph{el marinero}
  \item \emph{Recall:} say the Spanish for the sailing boat, then the Spanish for to row, then say \emph{el marinero} again
\end{itemize}

\begin{tcolorbox}[breakable,colback=teal!4,colframe=teal!35!black,title={Before you move on}]
What does el marinero mean? (\textquotedblleft{}The sailor\textquotedblright{}.) Say it once more.
\end{tcolorbox}
\section[el capitán]{el capitán — the captain}

\label{lesson:ES-C659-capitan}



\begin{quote}
Before the new one: say the Spanish for to row, then the Spanish for the sailor.
\end{quote}

\subsection*{You'll want to know: el capitán}
\textbf{el capitán} — \textquotedblleft{}the captain\textquotedblright{}.

\begin{quote}
El capitán habla con los pasajeros.
\end{quote}

\begin{quote}
\emph{The captain is talking to the passengers.}
\end{quote}

Of a ship, a plane or a team.

\begin{cousinweb}
\emph{capitán} comes from Late Latin \emph{capitāneus} \textquotedblleft{}chief\textquotedblright{} (EtymDB), built on \emph{caput} \textquotedblleft{}head\textquotedblright{}.
\end{cousinweb}

\begin{grammarlens}[title={What you've built}]
One more word for telling, explaining and describing at B1.
\end{grammarlens}

\subsection*{Your turn}
\begin{itemize}
\raggedright
  \item \emph{Say it:} \emph{el capitán}
  \item \emph{Say it:} \emph{el capitán}, once more
  \item \emph{Say it:} a sentence of your own with \emph{el capitán}
  \item \emph{Recall:} say the Spanish for to row, then the Spanish for the sailor, then say \emph{el capitán} again
\end{itemize}

\begin{tcolorbox}[breakable,colback=teal!4,colframe=teal!35!black,title={Before you move on}]
What does el capitán mean? (\textquotedblleft{}The captain\textquotedblright{}.) Say it once more.
\end{tcolorbox}
